% SC1004 Linear Algebra. Print on one A4 sheet, landscape, double-sided.
\documentclass[10pt,a4paper,landscape]{article}

\usepackage[T1]{fontenc}
\usepackage{lmodern}
\usepackage[margin=7.5mm]{geometry}
\usepackage{amsmath,amssymb,array,xcolor}

\renewcommand{\familydefault}{\sfdefault}
\definecolor{ink}{HTML}{19252C}
\definecolor{accent}{HTML}{173E50}
\definecolor{rulegray}{HTML}{A7B7BF}
\definecolor{wash}{HTML}{EDF2F4}
\pagestyle{empty}
\setlength{\parindent}{0pt}
\setlength{\parskip}{1.9pt}
\setlength{\fboxsep}{4pt}
\setlength{\tabcolsep}{3pt}
\setlength{\emergencystretch}{1em}
\raggedbottom

\newlength{\sheetcol}
\setlength{\sheetcol}{\dimexpr(\textwidth-8mm)/3\relax}
\newcommand{\Col}{\operatorname{Col}}
\newcommand{\Row}{\operatorname{Row}}
\newcommand{\Null}{\operatorname{N}}
\newcommand{\RREF}{\operatorname{rref}}
\newcommand{\rank}{\operatorname{rank}}
\newcommand{\spn}{\operatorname{span}}
\newcommand{\proj}{\operatorname{proj}}
\newcommand{\tr}{\operatorname{tr}}
\newcommand{\diag}{\operatorname{diag}}
\newcommand{\RR}{\mathbb{R}}
\newcommand{\CC}{\mathbb{C}}
\newcommand{\norm}[1]{\lVert #1\rVert}
\newcommand{\abs}[1]{\lvert #1\rvert}
\newcommand{\smat}[1]{\bigl[\begin{smallmatrix}#1\end{smallmatrix}\bigr]}
\newcommand{\formula}[1]{\par\noindent\(\displaystyle #1\)\par}
\newcommand{\topic}[1]{%
  \par\addvspace{4.0pt}{\color{accent}\bfseries #1\par}\nobreak\vspace{-1pt}}
\newcommand{\labelside}[2]{%
  \noindent{\color{accent}\bfseries\fontsize{17}{19}\selectfont SC1004 Linear Algebra}%
  \hfill{\color{accent}\bfseries\fontsize{9}{11}\selectfont #1\enspace\textbar\enspace #2}\par
  \vspace{1.5pt}{\color{rulegray}\hrule height .55pt}\vspace{2.3pt}}
\newcommand{\route}[1]{%
  \noindent\colorbox{wash}{\parbox{\dimexpr\textwidth-2\fboxsep\relax}{#1}}\par\vspace{3.8pt}}
\newenvironment{panel}[1]{%
  \begin{minipage}[t]{\sheetcol}\vspace{0pt}%
  \raggedright\setlength{\parskip}{1.9pt}%
  {\color{accent}\bfseries\fontsize{11}{12.5}\selectfont \strut #1\par}%
  \vspace{1.2pt}{\color{rulegray}\hrule height .45pt}\vspace{1.5pt}%
}{\end{minipage}}
\newcommand{\cue}[1]{\textbf{#1}}
\newcommand{\rowend}{\tabularnewline[1.4pt]}

% Use the standard fonts on full TeX installations. A partial macOS
% XeTeX cache can use local STIX glyphs for the missing math fonts.
\newif\ifpartialfontcache
\ifdefined\XeTeXversion
  \IfFileExists{lmex10.pfb}{}{%
    \IfFileExists{/System/Library/Fonts/Supplemental/STIXTwoMath.otf}{%
      \partialfontcachetrue
    }{}%
  }
\fi
\ifpartialfontcache
\makeatletter
\input{ot1lmr.fd}
\input{omllmm.fd}
\input{omslmsy.fd}
\makeatother
\DeclareFontShape{OT1}{lmr}{m}{n}{<->rm-lmr9}{}
\DeclareFontShape{OML}{lmm}{m}{it}{<-6.5>lmmi5 <6.5-10>lmmi9 <10->lmmi10}{}
\DeclareFontShape{OMS}{lmsy}{m}{n}{<-6.5>lmsy5 <6.5-10>lmsy8 <10->lmsy10}{}
\font\sheetSTIX="STIX Two Math" at 8pt
\font\sheetSTIXDisplay="STIX Two Math" at 10.5pt
\font\sheetSTIXBrackets="STIX Two Math" at 15pt
\renewcommand{\RR}{\mathord{\hbox{\sheetSTIX\char"211D}}}
\renewcommand{\CC}{\mathord{\hbox{\sheetSTIX\char"2102}}}
\renewcommand{\sum}{\mathop{\mathchoice{\vcenter{\hbox{\sheetSTIXDisplay\char"2211}}}{\vcenter{\hbox{\sheetSTIX\char"2211}}}{\hbox{\sheetSTIX\char"2211}}{\hbox{\sheetSTIX\char"2211}}}\displaylimits}
\renewcommand{\prod}{\mathop{\mathchoice{\vcenter{\hbox{\sheetSTIXDisplay\char"220F}}}{\vcenter{\hbox{\sheetSTIX\char"220F}}}{\hbox{\sheetSTIX\char"220F}}{\hbox{\sheetSTIX\char"220F}}}\displaylimits}
\renewcommand{\smat}[1]{\mathopen{\vcenter{\hbox{\sheetSTIXBrackets[}}}\begin{smallmatrix}#1\end{smallmatrix}\mathclose{\vcenter{\hbox{\sheetSTIXBrackets]}}}}
\fi

\begin{document}
\fontsize{8.65}{10.25}\selectfont
\color{ink}
\labelside{FRONT}{Systems, determinants and coordinates}
\route{\cue{Start with what the question gives you.}\quad
  Pivots $\to$ solutions and dimensions \quad
  Determinant $\to$ invertibility and scale \quad
  Basis $\to$ coordinates and matrix of a map}

\noindent\begin{panel}{Rank, pivots and solutions}
\topic{What the equation says}
For $A=[a_1\,\cdots\,a_n]\in\RR^{m\times n}$,
$Ax=\sum_i x_i a_i$ is a combination of the columns.
The unknowns $x_i$ are the weights. To test whether $b$ is
in their span, solve $Ax=b$ by reducing $[A\mid b]$.

\topic{What each pivot tells you}
In REF, zero rows come last and leading entries move right
as you go down. RREF is unique. It has pivot $1$s and zeros elsewhere
in each pivot column. Swap rows, scale by a nonzero scalar,
or add a multiple of another row. Apply each operation to $b$ too.

\begin{tabular}{@{}>{\raggedright\arraybackslash}p{.43\linewidth}>{\raggedright\arraybackslash}p{.54\linewidth}@{}}
\cue{What you find} & \cue{What you can conclude}\rowend
$[0\ \cdots\ 0\mid c]$, $c\ne0$ & No exact solution.\rowend
Consistent, $r=n$ & Unique solution. Back-substitute.\rowend
Consistent, $r<n$ & Infinitely many. Set $n-r$ variables free.\rowend
Each column of $A$ has a pivot & Columns are LI. The map is one-to-one.\rowend
Each row of $A$ has a pivot & Columns span $\RR^m$. The map is onto.\rowend
\end{tabular}

Here $r=\rank A$ counts pivots of $A$. For a particular $b$,
\formula{Ax=b\text{ is consistent}\iff b\in\Col A\iff\rank[A\mid b]=r.}
Find one solution $x_p$. All solutions are $x_p+\Null(A)$.
When $b=0$, the solution set is the null space itself.

\topic{Get bases and dimensions from rank}
\begin{tabular}{@{}p{.27\linewidth}p{.11\linewidth}>{\raggedright\arraybackslash}p{.56\linewidth}@{}}
\cue{Space} & \cue{Dim.} & \cue{Basis recipe}\rowend
$\Col A\subseteq\RR^m$ & $r$ & Pivot columns of the \textbf{original} $A$. \rowend
$\Row A\subseteq\RR^n$ & $r$ & Nonzero rows of REF. \rowend
$\Null(A)\subseteq\RR^n$ & $n-r$ & Solve $Ax=0$. Set each free variable to $1$ with the others $0$ in turn. \rowend
$\Null(A^{\top})\subseteq\RR^m$ & $m-r$ & Solve $A^{\top}y=0$ the same way. \rowend
\end{tabular}

$\Row A=\Col(A^{\top})$ and $\rank A^{\top}=r$.
Row operations preserve the row space, null space and relations
among columns. \cue{Rank--nullity} is $r+\dim\Null(A)=n$.

\topic{Use the same rank test in other problems}
$\Null(A)=\{0\}$ gives unique coefficients for a consistent
system or a least-squares fit. A nonzero null space of
$A-\lambda I$ gives eigenvectors. If all you know is the shape,
$m<n$ forces free variables and $m>n$ forces a proper column space.
\end{panel}%
\hfill
\begin{panel}{Invertibility and determinants}
\topic{For square \(A\), these tests are equivalent}
\formula{\begin{aligned}
 & A\text{ invertible}\iff\det A\ne0\iff\rank A=n\\
 & \iff\RREF(A)=I\iff\Null(A)=\{0\}\\
 & \iff\text{columns (or rows) form a basis of }\RR^n\\
 & \iff 0\text{ is not an eigenvalue of }A\\
 & \iff Ax=b\text{ has a unique solution for every }b.
\end{aligned}}
Any one proves all the others. A left or right inverse alone
also proves invertibility for square $A$.

\cue{If \(\det A=0\)} the map collapses volume,
at least one input direction goes to $0$, and any attainable
$b$ has infinitely many preimages. The zero eigenspace is $\Null(A)$.

\topic{Choose a way to find the determinant}
\begin{tabular}{@{}>{\raggedright\arraybackslash}p{.27\linewidth}>{\raggedright\arraybackslash}p{.69\linewidth}@{}}
\cue{What you see} & \cue{Calculation}\rowend
$2\times2$ & $ad-bc$ for entries $a,b$ on top and $c,d$ below.\rowend
Triangular & Product of diagonal entries.\rowend
Sparse row $i$ & $\sum_j(-1)^{i+j}a_{ij}\det A_{ij}$. Delete row $i$ and column $j$ for $A_{ij}$.\rowend
Eigenvalues & $\det A=\prod_i\lambda_i$, counting multiplicity.\rowend
\end{tabular}

\cue{Row or column operations} A swap flips sign. Scaling
one row by $c$ multiplies the determinant by $c$. Adding a
multiple of another row preserves it. Track these changes
when reducing to triangular form. Equal or zero rows give $0$.

$\det(AB)=\det A\,\det B$, $\det A^{\top}=\det A$,
$\det(cA)=c^n\det A$ and $\det(A^{-1})=1/\det A$ when invertible.
$\abs{\det A}$ is the area or volume scale factor.
For a nonzero determinant, the sign gives orientation.

\topic{Use an inverse or reuse a factorisation}
\cue{Undo a map} Row reduce $[A\mid I]\to[I\mid A^{-1}]$,
then $x=A^{-1}b$. For $2\times2$,
$A^{-1}=\smat{d&-b\\-c&a}/\det A$ when $\det A\ne0$.
$(AB)^{-1}=B^{-1}A^{-1}$ and $(A^\top)^{-1}=(A^{-1})^\top$.

\cue{Apply a row operation} Apply it to $I$ to get $E$.
$EA$ performs the operation on $A$, and $E^{-1}$ undoes it.

\cue{Solve for many \(b\)} Factor $PA=LU$, where $P$ records
row swaps ($I$ if none), $L$ stores elimination multipliers
below its unit diagonal, and $U$ is the result of elimination.
Solve $Ly=Pb$ forward, then $Ux=y$ back.
\end{panel}%
\hfill
\begin{panel}{Bases, coordinates and linear maps}
\topic{Test a set or extract a basis}
A vector space has closed addition and scalar multiplication,
with the vector-space axioms. For a subset $W\subseteq V$
under the same operations, test $0\in W$ and closure under
$\alpha u+\beta v$ for all $u,v\in W$ and scalars.
Spans, kernels and ranges are subspaces.

\cue{Span} is the set of attainable combinations.
\cue{Linear independence (LI)} means $\sum_i c_iv_i=0$
forces every $c_i=0$. A nontrivial relation lets you remove
a dependent vector while preserving the span.

\cue{A basis} is independent and spans. Put vectors into
columns in known coordinates and keep the original pivot columns
to get a basis for the span. Its size is the dimension.
With exactly $\dim W$ vectors in $W$, either LI or spanning
establishes a basis. For polynomials of degree $\le d$,
use the coefficient basis $\{1,t,\ldots,t^d\}$.

\topic{Use a basis to describe the same vector}
For a basis $\mathcal B=(b_1,\dots,b_n)$ of $\RR^n$,
put $P_{\mathcal B}=[b_1\,\cdots\,b_n]$.
Synthesize $x=P_{\mathcal B}[x]_{\mathcal B}$.
Extract weights by solving that equation, or apply $P_{\mathcal B}^{-1}$.
An orthonormal basis gives weights directly via $P_{\mathcal B}^{\top}x$.
\formula{[x]_{\mathcal C}=P_{\mathcal C}^{-1}P_{\mathcal B}[x]_{\mathcal B}.}
Column $i$ of the coordinate-change matrix is $[b_i]_{\mathcal C}$.

\topic{Build or change the matrix of a map}
Test $T(\alpha u+\beta v)=\alpha T(u)+\beta T(v)$.
For a linear $T\colon\RR^n\to\RR^m$, $T(x)=Ax$ where
column $i$ of $A$ is $T(e_i)$.
$\ker T=\Null(A)$ identifies directions sent to zero.
$\operatorname{range}T=\Col A$ gives reachable outputs.
\formula{[T]_{\mathcal C\leftarrow\mathcal B}=P_{\mathcal C}^{-1}AP_{\mathcal B}.}
Read right to left. Convert input coordinates, apply $A$,
then convert output coordinates. With the same basis
at both ends this is similarity. An eigenbasis makes it diagonal.

\cue{Compose maps} $AB$ applies $B$ first. The inner dimensions
must agree, $(AB)_{ij}=\sum_k a_{ik}b_{kj}$ and column $j$ is $Ab_j$.
$(AB)^\top=B^\top A^\top$. Add matrices of the same size
entrywise. Transpose swaps rows and columns.

\topic{Read a line or plane as an affine space}
Use $x=p+tv$ for a line, $x=p+su+tv$ for a plane with LI
directions, or $a^\top(x-p)=0$ for a hyperplane with normal $a\ne0$.
$u$ is parallel to $v\ne0$ when $u=cv$. Given a normal,
find directions by solving $a^\top v=0$.
\end{panel}

\newpage
\labelside{BACK}{Projections, eigendirections and frequencies}
\route{\cue{Choose a basis that fits the operation.}\quad
  Orthogonal $\to$ dot products give weights and closest fits \quad
  Eigenvectors $\to$ powers by scaling \quad
  Fourier $\to$ frequency analysis}

\noindent\begin{panel}{Orthogonality and closest fits}
\topic{Start with a dot product}
$u^\top v=\sum_i u_i v_i$. For nonzero vectors,
$\cos\theta=u^\top v/(\norm u\norm v)$.
Zero dot product means orthogonal. $\norm v=\sqrt{v^\top v}$,
$d(u,v)=\norm{u-v}$ and $v/\norm v$ is unit for $v\ne0$.
$\abs{u^\top v}\le\norm u\norm v$ (Cauchy--Schwarz) and
$\norm{u+v}\le\norm u+\norm v$.

\topic{Extract weights and project}
$W^\perp$ contains the vectors orthogonal to every vector in $W$.
In $\RR^n$, its dimension is $n-\dim W$.
$\Col(A)^\perp=\Null(A^\top)$ and $\Row(A)^\perp=\Null(A)$.

With an orthogonal basis $u_i$ of $W$, take the component
of $y$ along each basis vector and add the components.
\formula{c_i=\frac{u_i^\top y}{u_i^\top u_i},\qquad
  \hat y=\proj_W y=\textstyle\sum_i c_i u_i.}
$y=\hat y+e$ with $e\in W^\perp$. Pythagoras gives
$\norm{y-w}^2=\norm e^2+\norm{\hat y-w}^2$ for $w\in W$.
So $\hat y$ is the unique closest point and $\norm e$ is the
distance. For $p+W$, project $y-p$ and add $p$.

\topic{Get an orthonormal basis and a QR factorisation}
Nonzero orthogonal vectors are LI. $Q^\top Q=I$ means
orthonormal columns, so $Q$ preserves dot products. $QQ^\top$
projects onto $\Col Q$. For square $Q$, $Q^{-1}=Q^\top$.

\cue{Gram--Schmidt} Subtract earlier components using
$v_j=a_j-\sum_{i<j}q_i(q_i^\top a_j)$, then $q_j=v_j/\norm{v_j}$.
Skip zero remainders when building a basis. For $n$ LI columns,
$A=QR$ with $Q$ of size $m\times n$ and $R=Q^\top A$
upper triangular. $r_{ij}=q_i^\top a_j$, $r_{jj}=\norm{v_j}$.
Full QR completes $Q$ to an orthonormal basis of $\RR^m$.
In either form $\rank A=\rank R$.

\topic{Find the closest attainable \(Ax\)}
\cue{Least squares} Minimises $\norm{b-Ax}^2$. It projects
$b$ onto $\Col A$. The residual is orthogonal to the columns.
\formula{e=b-A\hat x,\quad A^\top e=0
  \iff A^\top A\hat x=A^\top b.}
Solve by elimination. For LI columns, use
$\hat x=(A^\top A)^{-1}A^\top b$ or back-solve
$R\hat x=Q^\top b$ from reduced QR. For dependent columns,
all minimisers are $\hat x+\Null(A)$.
The fitted vector $A\hat x$ is unique in either case.
$e=0$ exactly when the original system is consistent.
For a model $\sum_j c_j\phi_j(t)$, set $A_{ij}=\phi_j(t_i)$
and $b_i=y_i$ to fit the coefficients $c_j$.
\end{panel}%
\hfill
\begin{panel}{Eigenvalues and matrix structure}
\topic{Turn an eigenvalue question into a null-space question}
$Av=\lambda v$ with $v\ne0$ means $A$ scales that direction.
\formula{\det(A-\lambda I)=0,\qquad E_\lambda=\Null(A-\lambda I).}
Roots give eigenvalues. Solve the shifted homogeneous system
to get eigenspace bases. For triangular $A$, read the diagonal.
For $2\times2$, use $\det(tI-A)=t^2-(\tr A)t+\det A$.

\cue{Algebraic multiplicity} counts copies of a root.
\cue{Geometric multiplicity} is $\dim E_\lambda=n-\rank(A-\lambda I)$.
It lies between $1$ and the algebraic multiplicity.
Eigenvectors from distinct eigenvalues are LI.

\topic{Decide whether you can diagonalise}
Need $n$ independent eigenvectors over the chosen field.
When the characteristic polynomial splits, check that the
two multiplicities match at every root. $n$ distinct eigenvalues
suffice. Set $P=[v_1\,\cdots\,v_n]$ and
$D=\diag(\lambda_1,\dots,\lambda_n)$ in matching order.
\formula{AP=PD\ \Longrightarrow\ A=PDP^{-1},\qquad A^k=PD^kP^{-1}.}
For $x_{k+1}=Ax_k$, turn $x_0$ into eigenbasis coordinates with
$P^{-1}$, then scale coordinate $i$ by $\lambda_i^k$.
Active modes decay for $\abs\lambda<1$, keep their magnitude
for $\abs\lambda=1$ and grow for $\abs\lambda>1$.

\cue{Read invariants} Similar matrices share the characteristic
polynomial, determinant, trace and rank. $\tr A=\sum_i\lambda_i$.
$A^\top$ shares the spectrum. The same eigenvector has
eigenvalue $\lambda^k$ for $A^k$, $\lambda+c$ for $A+cI$,
or $1/\lambda$ for invertible $A^{-1}$.

\topic{Let structure tell you what is possible}
\cue{Real symmetric \(S\)} satisfies $S^\top=S$. It has real
eigenvalues and an orthonormal eigenbasis, so $S=QDQ^\top$. Positive definite
means $x^\top Sx>0$ for $x\ne0$, equivalent to all eigenvalues
being positive for symmetric $S$.

\cue{Gram matrix} $G=A^\top A$ gives $x^\top Gx=\norm{Ax}^2$.
$\Null(G)=\Null(A)$, so $G$ is invertible (and positive definite)
exactly when the columns are LI. Use this to test uniqueness of
the least-squares coefficients.

\cue{Orthogonal projection} $P^\top=P$ and $P^2=P$ together
characterise an orthogonal projector. $P$ fixes its range
and sends its orthogonal complement to zero.
Use $P=QQ^\top$ for an orthonormal basis of the target.
For LI columns of $A$, $P=A(A^\top A)^{-1}A^\top$ projects
onto $\Col A$. Then $\hat b=Pb$ and $e=(I-P)b$.

\cue{Complex pair for real \(A\)} Conjugate eigenvalues and
eigenvectors pair up. A non-real $\lambda=a+jb$ acts on the
real plane with basis $(\Re v,-\Im v)$ by the block
$\smat{a&-b\\b&a}$, scale $\abs\lambda$ and rotation $\arg\lambda$.
\end{panel}%
\hfill
\begin{panel}{Complex numbers and frequency}
\topic{Choose a form for the operation}
$z=a+jb=re^{j\theta}$, $j^2=-1$, $\bar z=a-jb$,
$r=\sqrt{a^2+b^2}$ and $\theta=\operatorname{atan2}(b,a)$ for $z\ne0$.
Euler gives $e^{j\theta}=\cos\theta+j\sin\theta$.
Add componentwise. Multiply by multiplying moduli and
adding angles. Divide by dividing moduli and subtracting
angles, or use $1/z=\bar z/r^2$.
\formula{z^p=r^pe^{jp\theta},\qquad
 \text{\(p\)th roots }\;r^{1/p}e^{j(\theta+2\pi\ell)/p},\quad
 \ell=0,\dots,p-1.}

\topic{Apply orthogonality with conjugation}
$u^H=\bar u^\top$ and $\langle u,v\rangle=u^Hv$.
$\norm u^2=u^Hu=\sum_n\abs{u_n}^2$.
The coefficient along an orthogonal basis vector $u$ is
$u^Hy/(u^Hu)$. For an orthonormal basis, the denominator is $1$.

\topic{DFT gives weights. IDFT rebuilds samples.}
Let $\Omega_k=2\pi k/N$ and $w_k[n]=e^{j\Omega_k n}$,
for $n,k=0,\dots,N-1$. $w_k^Hw_\ell=N\delta_{k\ell}$,
so this is an orthogonal basis and $\norm{w_k}^2=N$.
\formula{X[k]=\sum_{n=0}^{N-1}x[n]e^{-j\Omega_k n}=w_k^Hx.}
\formula{x[n]=\frac1N\sum_{k=0}^{N-1}X[k]e^{j\Omega_k n}.}
The synthesis weight is $X[k]/N$. With $W_{kn}=e^{-j\Omega_k n}$,
$X=Wx$ and $W^{-1}=W^H/N$. Use row $k$ of $W$ for a
single bin or the corresponding column of $W^H$ and $1/N$
to reconstruct that bin's contribution.

\topic{Read the signal from the bins}
\cue{Frequency and direction} Reduce the phase step modulo
$2\pi$ to $(-\pi,\pi]$. The synthesis basis uses $+\Omega_k$.
The analysis row uses $-\Omega_k$. A step strictly between
$0$ and $\pi$ rotates anticlockwise. A step strictly between
$-\pi$ and $0$ rotates clockwise. A zero step is constant
and $\pi$ alternates. Bin $k>N/2$ has signed index $k-N$
in synthesis. Bins $k$ and $k+N$ are aliases.

\cue{Amplitude and phase} For a complex tone $a\,w_k$,
only bin $k$ is active, with $X[k]=Na$. Read amplitude
$\abs{X[k]}/N$ and phase $\arg X[k]$. For real inputs,
$X[N-k]=\overline{X[k]}$ modulo $N$. A paired real sinusoid
has amplitude $2\abs{X[k]}/N$. At DC and Nyquist use the
signed weight $X[k]/N$.

\cue{Check a result} $X[0]/N$ is the sample mean.
Parseval gives $\sum_n\abs{x[n]}^2=\sum_k\abs{X[k]}^2/N$.
\end{panel}
\end{document}
